\sscp{Issues with glosses}{02-04-04}
Some elusive terms or glosses are central to debates about Austronesian prehistory.
For example, \citet{Fox2011} points out that an 1889 Dutch source
for certain \ili{Ngaju} Dayak sibling terms has not been able to be
verified by other anthropologists (James Fox pers. comm. October 2020),
nor is it known by native speakers.
Verifying the term (or not) is important since it appears to
be an anomaly in the distribution of cognate forms,
most of which are found far to the \ili{east} in Linguistic Wallacea.
A report from nearby Uut Danum shows possible cognates restricted
to formal speech and oral epics.
The terms in Borneo are elusive at best,
and may have been present in certain language(s) in the past,
but are no longer current.
Yet the existence and accurate glossing of such sibling terms
denoting relative age and opposite sex are
critical to understanding early Austronesian
social structure and transformations
in different regions (see \citealp{Fox2015}).

Difficulties with glosses impinge on precisely identifying \it{flora
and fauna} for comparative purposes.
For instance, word lists with ‘grass’ often list the first answer
given, but the investigator may or may not explore further to see if
it is a generic term,
or if there are different terms for saw-toothed vs.
smooth edged varieties -- which often have different historical etyma.
Is the word given for ‘bamboo’ a generic term,
or does it refer to a certain species of bamboo? Word lists with
‘uncle’ often do not distinguish between ‘mother’s brother’
and ‘father’s brother’, which are not distinguished in the languages
of elicitation, but are consistently distinguished in the social structures
and lexicons of our target region.
Investigators taking word lists may assume sibling terms have
only ‘elder sibling’ and ‘younger sibling’ distinctions
and fail to check if there are any ‘opposite sex sibling’,
‘older sibling (same sex)’,
or ‘younger sibling (same sex)’ terms which may not exist where
the investigator is from or in the languages of elicitation,
but are common in the target region.

PMP *duRi is reconstructed meaning ‘thorn’,
which is the overwhelming gloss in languages in the northern
and western Malayo-Polynesian regions.
East of the Blust line the dominant gloss shifts to ‘bone’
and continues out into the Pacific (see \frf{15-02}).
Yet some investigators in eastern Indonesia who have Austronesian
awareness about the etymon meaning ‘thorn’,
have glossed reflexes of *duRi only as ‘fish bone’,
whereas investigators without such bias or awareness gloss them
simply as ‘bone’ in the same language.
It is the word given when working through
lists of human body parts or pig anatomy.

\ili{As} mentioned before, our target region sits between two biogeographical
zones, and also has flora
and fauna endemic to the region.
While there are resources on the flora
and fauna of the New Guinea region (Sahulland),
and of South\ili{east} Asia (Sundaland),
there are few focused on Wallacea.

In discussing discrepancies between the distribution of bandicoots
(the rodents) and the distribution of words glossed as ‘bandicoot’,
\citet[262--263]{Schapper2011} observes, “Six languages with claimed
reflexes in this region do not occur in the area where bandicoots are found.
Furthermore, the terms do not denote ‘bandicoot’.
\citet[110]{Blust2002} makes light of the discrepancy between the meaning
of reflexes and that of his reconstruction,
stating that ‘in some cases reflexes of this form are said to
indicate the phalanger, or cuscus, and it is
not entirely clear whether this is due to
semantic change or to confusion on the part of the lexicographer.’
In fact, with only a single exception,
all reflexes of *mans(aə)r ‘bandicoot’ west of New Guinea denote
‘cuscus’ and not ‘bandicoot’, even in the area where bandicoots are found.”

So Blust appeals to “confusion on the part of the lexicographer”,
yet he also appeals to the opposite when a critical piece of
data supporting his hypothesis is challenged.
\citet[263]{Schapper2011} points out that the \ili{Leti}-\ili{Moa} term \it{mada
{\tl} made} has been incorrectly interpreted as ‘bandicoot’ by
Blust and the correct gloss as ‘cuscus’ has been verified by
\citet{vanEngelenhoven2004}
and that the \it{madar} term reported for \ili{Damer} whose gloss is
translated as ‘marsupial sp.,
bandicoot’ is actually outside the area where bandicoots are found.

\citeauthor{Blust2012b}'s (\citeyear[264]{Blust2012b}) response to this is curious.
Ignoring the information that \citeauthor{vanEngelenhoven2004} verified the \ili{Leti}-\ili{Moa}
gloss as ‘cuscus’, he says, “The one clear discrepancy that remains in this table
is the gloss for \ili{Leti}-\ili{Moa} \it{mada/made},
which \citet{Jonker1932} cites as buidelrat (‘pouch rat’ = bandicoot).
{\ldots} It is odd that Jonker,
who was a very experienced lexicographer
and a scholar well acquainted with the land,
peoples, and languages of the Lesser Sundas,
would commit such an error.” We note that Jonker’s dictionaries
(like most first dictionaries of any language) are full of dozens
of vague entries saying “soort X [soort boom]” ‘a kind of X [a kind of tree]’.
That does not diminish his skill as a lexicographer,
it only points out he was working with many limitations in the
specific identification of flora and fauna.

The debate between \citeauthor{Blust2012b}
and \citeauthor{Schapper2011} illustrates another common methodological problem
of filtering obscure glosses through multiple languages,
and the caution and grace needed in sorting things out.
For example, were the individual Holle word lists done from the
late 1800s to early 1900s taken using Dutch as the language of
elicitation, or \ili{Malay}? Or did some use Dutch
and some use \ili{Malay}? How proficient was the investigator taking
the word list (Dutch investigators using \ili{Malay}; local investigators
using Dutch)? What about proficiency in the regional variety
of \ili{Malay}, which can have significant differences in the vocabulary from
Standard \ili{Malay}? (See examples below.) How proficient in the language
of elicitation (Dutch or \ili{Malay}) was the informant supplying the
data for the word list (some of whom were second language speakers
of the target language).
And then how well do \ili{English} speakers using the materials handle
Dutch or \ili{Malay} glosses from a previous era?

For example, Dutch \it{buidelrat} parses literally to ‘pouch
rat’, so assuming it means some type of ‘marsupial’ is reasonable.
And Dutch \it{rat} means ‘rat’,
so assuming ‘pouch rat’ refers to ‘bandicoot’ is not unreasonable.
But several modern Dutch dictionaries give ‘possum’ (or ‘opossum’)
for \it{buidelrat}.\footnote{
		For example, putting \it{buidelrat} into the Dutch version of Wikipedia redirects
		to opossum (\it{Didelphidae}) -- the family of marsupials found
		in the Americas. (‘Opossums’, Wikipedia.
		nl.wikipedia.org/wiki/ Opossums. Accessed 22 June 2023)}
The Holle list master \citep[50]{Stokhof1980a} equates: item {\#}819
Dutch: \it{buideldier} = \ili{English}: \it{marsupial} = Indonesian: \it{kuskus} ‘cuscus’.
Which language was used for elicitation? The Dutch term could
cover any marsupial, whereas the \ili{Malay} is limited to the \it{Phalanger} ‘cuscus’.
\ili{Ambon} \ili{Malay} has \it{kusu} for ‘cuscus’,
perhaps to disambiguate it from \it{kusu-kusu} ‘kunai grass,
sword grass’ (Indonesian: \it{alaŋ-alaŋ}),
which in rapid speech is pronounced \it{kuskusu}
often with an utterance final decrescendo.
\citet[304f]{vanEkris1864} provides independent evidence via \ili{Malay}
that Dutch \it{buidelrat} is directly associated with ‘cuscus’
during the era when many Dutch sources recorded languages for
which we have no other data.
For the \ili{Kamarian} language of southwest Seram (now said to be
extinct), \citeauthor{vanEkris1864} gives \it{maker} ‘buidelrat’,
\it{make mauriat} ‘gevlekte buidelrat [Ml.
kusu potar]’ (Ml. = \ili{Malay}), \it{make moon} ‘groote buidelrat’.
\ili{As} mentioned above \it{kusu} is the \ili{Ambon} \ili{Malay} term for ‘cuscus’.
And Dutch \it{gevlekte} = ‘spotted’,
so \it{make mauriat} ‘gevlekte buidelrat’ being explicitly identified
in \ili{Malay} as a type of \it{kusu} refers to the ‘spotted cuscus
(\it{Spilocuscus maculatus}; see \frf{02-26})’.
This is fairly clear evidence that the term \it{buidelrat} in
19th century Dutch sources can refer to ‘cuscus’.

\begin{figure}
	\centering
	\includegraphics[width=0.45\columnwidth]{02-26.jpg}
	\fcp{Spotted cuscus \it{(Spilocuscus maculatus)}}{02-26}
\end{figure}

So given the evidence from \citet{vanEkris1864}
and modern Dutch dictionaries,
both \citeauthor{Jonker1932} and \citeauthor{vanEngelenhoven2004} agree that \ili{Leti}
and \ili{Moa} \it{mada/made} refer to ‘cuscus’,
and we do not need to discuss whether Jonker made an error or not.
Similar cross-linguistic confusion can be found with terms referring
to civet cats and bandicoots as well as types of birds,
bats, insects and plants.

To expand on the difficulties in identifying flora
and fauna, some might also characterise both Grimes
and Edwards as “experienced lexicographer[s]” (see Bibliography).
Yet we are happy to acknowledge difficulties in identifying terms
in the field for elusive animals such as the bandicoot
and civet cat, or even more common plants or birds or insects.
There are many variables.
For example: coastal people tend to know very little about the
flora and fauna of the interior or the highlands,
and \it{vice versa} with mountain people,
especially on the larger islands.

The lexicographer can be reasonably fluent or even very fluent
in the local language,
but be talking about an unknown thing,
or limited to verbal descriptions.
For example, \ili{Buru} men in the interior have described to Grimes
a \it{faf bile} that is said to be distinct from long-legged
wild pigs (\it{faf aba}),
distinct from domestic pigs with slightly shorter legs (\it{faf
fena}), distinct from domestic pigs that have gone feral (\it{fafu melaban}),
distinct from wild pigs that have been caught young
and domesticated (\it{fafu emlahat}),
and distinct from the babirusa (\it{faf donit}).
The \it{faf bile} is said to live in the jungle
and “has short legs and fine hair”.
Several conversations with additional hunters who also spend
considerable time in the jungle have not added any new information.
More research is clearly needed.
But when is that likely to happen?

More commonly the comparative researcher has limited fluency
or may even be very fluent in a language of wider communication,
or an official language,
both of which add additional filters
and limitations in both directions.
Many people in remote communities often have no reason to talk
about local flora and fauna in the national language
and often don’t actually know the official language terms.
People in \ili{Buru} know a bit of \ili{Ambon} \ili{Malay}.
People in west Timor know a bit of \ili{Kupang} \ili{Malay}.
The \ili{Ambon} terms and the \ili{Kupang} terms for some of the same types
of plants and insects and fish can be quite different from each other due
to different patterns of contact.
Add to that the problem that most of the standard Indonesian
terms for flora and fauna in major or official dictionaries are biased towards
terms used in the Sumatra-Java region,
and are often oblivious to terms prevalent
in the different regions of eastern Indonesia.

For example, the Indonesian Bible term for ‘cow’ used in Genesis
41:1-4 for pharoah’s dream is \it{ləmbu},
which is not correctly identified as any type of ‘cow’
and sometimes not even as an animal by 70{\%} of 270 tertiary
level theology students and ministers tested in eastern Indonesia by Grimes.
A few recognised it as “an animal used in sacrifices”.
These are not uneducated people,
and they all know the Indonesian word \it{sapi} ‘cow’,
but they live in a region where nobody
uses \it{ləmbu} outside the Indonesian Bible.

Similarly in discussing ways to translate the story of the flood,
the standard Indonesian term for the ‘raven’ released by Noah
in Genesis 8:6--7 is \it{burung gagak}.
In several different remote language communities in eastern Indonesia,
when people were asked to describe what that might be,
there were several similar responses separated by time
and distance and language region saying it might be white or
speckled, and it eats fruit.
Again, these people are not ignorant.
They just don’t know that official language term.

When Grimes was assisting the Ministry of Education with multilingual
education materials in the \ili{Kemak} \ili{Marobo} area of Timor-Leste,
several farmers of both sexes took four days to come up with
the \ili{Kemak} word for ‘ladybug’ for a children’s book selected by
decision-makers in the city.
When they finally recalled it,
it was not a word recognised by most speakers.
A ladybug has no local cultural value
and figures in no folktales.
Several knew the Portuguese loan \it{joaninho} from school books,
because the insect has cultural value in Europe for children’s stories.

But even with the aid of language fluency
and drawings and digital photos the issue is not always clear.
People are used to seeing things in the jungle in colour,
in three dimensions, in their habitats, and in motion.
We are often asking them to interpret something static in two
dimensions out of context,
or from verbal description alone.
Line drawings without colour can also have some drawbacks.
When Grimes showed the black
and white drawing to five men who live in the remote
mountains and spend considerable time hunting,
three different names of doves/pigeons were identified,
discussed and argued over, and a fourth
possibility was named and then discarded.
It was not until he showed the colour version of the drawing (\frf{02-27})
that they eventually reached consensus that it was the \it{man rahe}.
But these visual aids to identifying local flora
and fauna are often unavailable,
or from different biogeographical zones,
or one may be limited to working verbally.


\begin{figure}
	\centering
	\includegraphics[width=0.6\columnwidth]{02-27.jpg}
	\fcp{Emerald dove \it{(Chalcophaps indica)}}{02-27}
\end{figure}

Grimes has also gotten different responses to questions like
“Describe the \it{goli} bird?” Some more experienced hunters
described it as a medium sized meat-eating bird with black wings
that hunts live prey and whose call is \it{ekrikik}.
A more home-body type described it as a seed-and-fruit eating
bird whose call is \it{wio}.
The confidence level is higher for the first one
and whose description probably points to the ‘black-winged
kite’ (\it{Elanus caeruleus hypoleucus Gould}).
But discrepancies mean that further checking is warranted,
and it is unknown when that might happen.

%\begin{figure}
    %\subfigure[]{
		%\includegraphics[width=\textwidth]{02-27a.jpg}
		%}
    %\subfigure[]{
		%\includegraphics[width=\textwidth]{02-27b.jpg}
		%}
	%\fcp{Emerald dove \it{(Chalcophaps indica)}}{02-27}
%\end{figure}

Having experienced many uncertainties in identifying flora
and fauna, and having gotten a few things wrong ourselves,
we are willing to countenance the possibility that some published
sources might not have gotten everything quite right,
particularly when it is anomalous to the rest of the data.
Subgrouping arguments should not hinge on disputed data,
and particularly not on disputed lexical data or their glosses.
